\problem{2.1}{Reality, parity, and stationary states}
\textit{Paraphrase.} Establish three properties of stationary solutions for real
one-dimensional, time-independent potentials.
For a separated solution $\Psi=\psi(x)\ee^{-\ii Et/\hbar}$, let
$E=E_R+\ii E_I$. Then
\begin{equation}
\|\Psi(t)\|^2=\ee^{2E_It/\hbar}\|\psi\|^2.
\end{equation}
A normalizable closed-system state preserves its norm, so $E_I=0$:
the energy is real. For real $V$, complex conjugation of
$-\hbar^2\psi''/(2m)+V\psi=E\psi$ shows that $\psi^*$ solves the same
equation. Hence $\Re\psi=(\psi+\psi^*)/2$ and
$\Im\psi=(\psi-\psi^*)/(2\ii)$ provide real solutions (at least one
nonzero) spanning the same real solution space. If $V(-x)=V(x)$, the
reflected function $\psi(-x)$ also solves the same equation.
Then
\begin{equation}
\psi_e(x)=\frac{\psi(x)+\psi(-x)}2,\qquad
\psi_o(x)=\frac{\psi(x)-\psi(-x)}2
\end{equation}
are even and odd eigenfunctions with the same energy.
\answer{Bound-state eigenfunctions can be chosen real and of definite
parity, and their energies are real.}
\check{In higher-dimensional or degenerate problems not every complex
eigenvector is a single parity eigenfunction; one chooses a suitable
basis within the eigenspace.}

\problem{2.2}{Why bound-state energy exceeds the potential minimum}
\textit{Paraphrase.} Prove a lower bound for any nontrivial normalizable eigenstate.
Multiply $\hat H\psi=E\psi$ by $\psi^*$ and integrate. Integration by
parts under boundary conditions giving vanishing surface terms gives
\begin{equation}
E\int|\psi|^2\dd x=
\frac{\hbar^2}{2m}\int|\psi'|^2\dd x+
\int V|\psi|^2\dd x
\ge V_{\min}\int|\psi|^2\dd x .
\end{equation}
For a normalized state,
\begin{equation}
E-V_{\min}=\frac{\hbar^2}{2m}\int|\psi'|^2\dd x
+\int(V-V_{\min})|\psi|^2\dd x.
\end{equation}
Both terms are nonnegative. Equality would require $\psi'$ to
vanish almost everywhere and the support to lie where $V=V_{\min}$;
a nonzero constant is not square-integrable on the real line.
\answer{$\boxed{E>V_{\min}}$ for a nontrivial normalizable bound
state on the infinite line under standard assumptions.}
Classically $E=T+V\ge V_{\min}$; the strict quantum inequality reflects
the impossibility of a normalizable state with everywhere vanishing
kinetic expectation.

\problem{2.3}{Excluding nonpositive well energies}
In the well $V=0$ for $0<x<a$ and $\psi(0)=\psi(a)=0$.
For $E=0$ the equation is $\psi''=0$, hence
$\psi=Ax+B$. The endpoint values yield $B=0$, then $A=0$.
For $E<0$, set $\kappa=\sqrt{-2mE}/\hbar>0$. Then
\begin{equation}
\psi=A\sinh(\kappa x)+B\cosh(\kappa x);
\quad \psi(0)=0\Rightarrow B=0;\quad
\psi(a)=0\Rightarrow A\sinh(\kappa a)=0.
\end{equation}
Since $\sinh(\kappa a)\ne0$, $A=0$ as well. Only the trivial state
remains. Thus $E_n=n^2\pi^2\hbar^2/(2ma^2)>0$.

\problem{2.4}{Position and momentum fluctuations in a box}
For $\phi_n=\sqrt{2/a}\sin(n\pi x/a)$,
\begin{align}
\avg{x}&=\frac2a\int_0^a x\sin^2(n\pi x/a)\dd x=\frac a2,\\
\avg{x^2}&=\frac2a\int_0^a x^2\sin^2(n\pi x/a)\dd x
=a^2\left(\frac13-\frac{1}{2n^2\pi^2}\right),\\
\avg p&=-\ii\hbar\int_0^a\phi_n\phi_n'\dd x
=-\frac{\ii\hbar}{2}[\phi_n^2]_0^a=0,\\
\avg{p^2}&=-\hbar^2\int_0^a\phi_n\phi_n''\dd x
=\frac{n^2\pi^2\hbar^2}{a^2}.
\end{align}
Therefore
\begin{equation}
\boxed{\sigma_x=a\sqrt{\frac1{12}-\frac1{2n^2\pi^2}},
\quad \sigma_p=\frac{n\pi\hbar}{a},\quad
\sigma_x\sigma_p=\hbar\sqrt{\frac{n^2\pi^2}{12}-\frac12}.}
\end{equation}
The expression increases with integer $n\ge1$, so the
ground state is closest to the Heisenberg bound:
$\sigma_x\sigma_p=\hbar\sqrt{\pi^2/12-1/2}\approx0.568\hbar$.
The strict inequality follows from the sine shape not being Gaussian.

\problem{2.5}{Equal two-state superposition in an infinite well}
Let $\Omega=\pi^2\hbar/(2ma^2)$ and
$\Psi(x,0)=A[\phi_1(x)+\phi_2(x)]$.
Orthonormality yields $1=|A|^2(1+1)$; choose $A=1/\sqrt2$.
The evolved state and its density are
\begin{align}
\Psi(x,t)&=\frac1{\sqrt2}
\left[\phi_1(x)\ee^{-\ii\Omega t}+
\phi_2(x)\ee^{-4\ii\Omega t}\right],\\
|\Psi|^2&=\frac12(\phi_1^2+\phi_2^2)
+\phi_1\phi_2\cos(3\Omega t).
\end{align}
The diagonal position elements are $a/2$, while
\begin{equation}
X_{12}=\int_0^a\phi_1(x)x\phi_2(x)\dd x
=\frac1a\int_0^a x[\cos(\pi x/a)-\cos(3\pi x/a)]\dd x
=-\frac{16a}{9\pi^2}.
\end{equation}
Hence
\begin{align}
\boxed{\avg{x}(t)=\frac a2-\frac{16a}{9\pi^2}\cos(3\Omega t)},\\
\boxed{\avg p(t)=m\frac{\dd\avg x}{\dd t}
=\frac{16ma\Omega}{3\pi^2}\sin(3\Omega t)}.
\end{align}
The oscillation amplitude is $16a/(9\pi^2)<a/2$ and angular
frequency $3\Omega$. Only $E_1=\hbar\Omega$ and
$E_2=4\hbar\Omega$ can be measured, with probability $1/2$
each; $\avg H=(E_1+E_2)/2=5\hbar\Omega/2$.

\problem{2.6}{Relative phase controls spatial oscillation}
Change the initial state to
$\Psi(x,0)=[\phi_1(x)+\ee^{\ii\varphi}\phi_2(x)]/\sqrt2$.
Time evolution gives
\begin{equation}
\Psi(x,t)=\frac{\phi_1\ee^{-\ii\Omega t}
+\ee^{\ii\varphi}\phi_2\ee^{-4\ii\Omega t}}{\sqrt2},
\end{equation}
so
\begin{align}
|\Psi|^2&=\frac12(\phi_1^2+\phi_2^2)
+\phi_1\phi_2\cos(3\Omega t-\varphi),\\
\boxed{\avg x(t)&=\frac a2-\frac{16a}{9\pi^2}
\cos(3\Omega t-\varphi).}
\end{align}
For $\varphi=\pi/2$ the oscillation is shifted by a quarter
period and $\avg x(0)=a/2$; for $\varphi=\pi$, the cosine term changes
sign and the initial spatial asymmetry reverses. The energy
probabilities remain $1/2$ and $1/2$. A global phase would cancel;
the \emph{relative} phase does not.

\problem{2.7}{Triangular initial state and spectral decomposition}
Let $\Psi(x,0)=A\min(x,a-x)$ on $[0,a]$, zero otherwise.
By symmetry,
\begin{equation}
1=2A^2\int_0^{a/2}x^2\dd x
=A^2\frac{a^3}{12}
\quad\Longrightarrow\quad A=\sqrt{\frac{12}{a^3}}.
\end{equation}
It has a peak $Aa/2$ at the midpoint. Its expansion coefficients are
\begin{align}
c_n&=\sqrt{\frac2a}A\left[
\int_0^{a/2}x\sin\frac{n\pi x}{a}\dd x+
\int_{a/2}^{a}(a-x)\sin\frac{n\pi x}{a}\dd x\right]\\
&=\boxed{\frac{4\sqrt6}{n^2\pi^2}
\sin\left(\frac{n\pi}{2}\right)}.
\end{align}
Thus all even terms vanish, and
\begin{equation}
\Psi(x,t)=\sum_{n=1,3,5,\ldots}c_n\phi_n(x)
\ee^{-\ii n^2\Omega t},\qquad
P(E_1)=|c_1|^2=\frac{96}{\pi^4}.
\end{equation}
Using $\sum_{n\ \mathrm{odd}\ge1}n^{-2}=\pi^2/8$,
\begin{align}
\avg H&=\sum_{n\ \mathrm{odd}}n^2\hbar\Omega|c_n|^2
=\hbar\Omega\frac{96}{\pi^4}\frac{\pi^2}{8}\\
&=\boxed{\frac{12\hbar\Omega}{\pi^2}
=\frac{6\hbar^2}{ma^2}}.
\end{align}
\check{The square-integrable jump in the first derivative does not
invalidate the spectral expansion; in Problem 2.39 it is treated
using distributions.}

\problem{2.8}{Energy measurement from a half-box constant state}
Here $\Psi(x,0)=A$ on $0<x<a/2$, and zero otherwise.
Normalization gives $|A|^2(a/2)=1$, so $A=\sqrt{2/a}$.
The ground-state overlap is
\begin{equation}
c_1=\sqrt{\frac2a}\sqrt{\frac2a}
\int_0^{a/2}\sin(\pi x/a)\dd x
=\frac{2}{a}\frac{a}{\pi}
=\frac2\pi.
\end{equation}
As the potential is time-independent, $|c_1|^2$ remains fixed:
\begin{equation}
\boxed{P\left(E=\frac{\pi^2\hbar^2}{2ma^2}\right)
=\frac4{\pi^2}\approx0.4053.}
\end{equation}

\problem{2.9}{Direct Hamiltonian expectation for a polynomial state}
For Example 2.2, $\Psi(x,0)=A x(a-x)$ with
$A=\sqrt{30/a^5}$. Inside the well,
\begin{equation}
\Psi''(x,0)=-2A,\qquad
\hat H\Psi=-\frac{\hbar^2}{2m}\Psi''
=\frac{\hbar^2 A}{m}.
\end{equation}
Then
\begin{align}
\avg H&=\int_0^a\Psi(x,0)\hat H\Psi(x,0)\dd x\\
&=\frac{\hbar^2A^2}{m}\int_0^a x(a-x)\dd x
=\frac{\hbar^2}{m}\frac{30}{a^5}\frac{a^3}{6}\\
&=\boxed{\frac{5\hbar^2}{ma^2}}.
\end{align}
This equals the spectral-sum result in Example 2.3.
